\begin{helper}
Fix the actual split-outside-blocker replacement surface for one step.  Thus
we have finite graphs \(G,G'\), activation--priority product measures
\(\nu,\nu'\) with the same parameter \(p=\gamma/\Delta\in[0,1]\), an
injective local embedding \(\phi:V\to V'\), a root \(r\), a finite candidate
set \(X\subseteq N_G(r)\), and a private-blocker map
\[
\beta:\{(x,z):x\in X,\ z\notin X\cup\{r\},\ xz\in E(G)\}\to V'
\]
with the usual replacement properties: \(\beta\) is injective, its values
avoid \(\phi(X\cup\{r\})\), \(\beta(x,z)\) is adjacent to \(\phi(x)\), no
other \(\phi(y)\) with \(y\in X\) is adjacent to it, \(\phi(r)\) is not
adjacent to it, and every \(G'\)-neighbor of \(\phi(x)\) is either an
embedded local neighbor or one of these private blockers.

Let \(z_0,\ldots,z_{n-1}\) be an injective list of outside blockers, and set
\[
  B_i=\{x\in X:xz_i\in E(G)\}.
\]
Assume each \(B_i\) is nonempty, each \(z_i\) lies in the source extension
coordinate layer, the source common and extension layers are disjoint and
contain \(X\cup\{r\}\) on the common side, and the target common layer is
\(\phi\) of the source common layer with a disjoint target extension layer
containing all private coordinates \(\beta(x,z_i)\) for \(x\in B_i\).  Let
\(C_\omega,C'_\omega\) be paired cells with common nonnegative weight
\(w(\omega)\), so
\[
  \operatorname{ofReal}(w(\omega))=\nu(C_\omega)
  =\nu'(C'_\omega).
\]

Now fix \(k<n\) and a cell index \(\omega\).  The source context
\(M^-_{k,\omega}\) is definitionally the fixed noncurrent source context
\[
\begin{aligned}
M^-_{k,\omega}=\{(A,\pi):&\ \exists x\in X,\ x\in A,\ 
   \pi(y)<\pi(x)\text{ for every }y\in (X\cup\{r\})\cap A
   \text{ with }xy\in E(G),\\
&\text{and for every }i>k\text{ with }x\in B_i,\ 
   z_i\text{ does not kill }x\}.
\end{aligned}
\]
Thus the only retained common blocker from the \(k\)-th handoff not included
in \(M^-_{k,\omega}\) is the current blocker \(z_k\).  The target context
\(M^+_{k,\omega}\) is definitionally its transported target analogue
\[
\begin{aligned}
M^+_{k,\omega}=\{(A',\pi'):&\ \exists x\in X,\ \phi(x)\in A',\
   \pi'(\phi(y))<\pi'(\phi(x))\text{ for every }y\in X\cup\{r\}
   \text{ with }\phi(y)\in A'\\
&\text{and }xy\in E(G),\text{ and for every }i<k\text{ with }x\in B_i,\ 
   \beta(x,z_i)\text{ does not kill }\phi(x)\}.
\end{aligned}
\]
The adjacent source event \(L_{k,\omega}\) is definitionally the event that
some \(x\in B_k\) is not killed by the single current common blocker \(z_k\).
The adjacent target event \(R_{k,\omega}\) is definitionally the event that
some \(x\in B_k\) is not killed by its independent private blocker
\(\beta(x,z_k)\).  These are statement-level definitions of the four displayed
sets, not proof-only interpretations.

Assume that \(B_k\) is represented by a nonempty finite label type
\(\alpha\), through a surjection \(a\mapsto x_a\) onto \(B_k\), and let
\(\pi_a\in[0,1]\) be the fixed priority of \(x_a\) after the noncurrent
context has been frozen.  Assume the current blocker coordinate, conditional
on the fixed source context, has a one-coordinate Bernoulli--uniform law
\(\eta\):
\[
\eta(\Omega)=1,\qquad
\eta\{a=1,\ q\in[0,1]\}=\operatorname{ofReal}(p),
\]
and
\[
\eta\{a=1,\ 0\le q\le t\}=\operatorname{ofReal}(pt)
\quad(0\le t\le1).
\]
Assume the target private blocker coordinates, conditional on the fixed
target context, have the independent product law \(\theta\): for every
activation set \(A\subseteq\alpha\) and every threshold vector
\(t_a\in[0,1]\) on \(A\),
\[
\theta\{a_b={\bf 1}_{b\in A}\ \forall b,\ 
        0\le q_b\le t_b\ \forall b\in A\}
=\operatorname{ofReal}\!\left(
  p^{|A|}(1-p)^{|\alpha|-|A|}\prod_{b\in A}t_b\right).
\]
Finally let \(P^-_{k,\omega},P^+_{k,\omega}\ge0\) be the normalized adjacent
source and target masses satisfying
\[
\operatorname{ofReal}(P^-_{k,\omega})
=\eta\{\neg(a=1,\ q\in[0,1],\ \pi_b<q\text{ for all }b\in\alpha)\},
\]
\[
\operatorname{ofReal}(P^+_{k,\omega})
=\theta\{\exists b\in\alpha:
  \neg(a_b=1,\ q_b\in[0,1],\ \pi_b<q_b)\},
\]
and whose unnormalized masses are the actual adjacent pieces on the same
cell surface,
\[
\operatorname{ofReal}(w(\omega)P^-_{k,\omega})
=\nu(L_{k,\omega}\cap M^-_{k,\omega}\cap C_\omega),
\]
\[
\operatorname{ofReal}(w(\omega)P^+_{k,\omega})
=\nu'(R_{k,\omega}\cap M^+_{k,\omega}\cap C'_\omega).
\]
Then
\[
  P^-_{k,\omega}\le P^+_{k,\omega}.
\]

This is a fixed-\(k\), fixed-\(\omega\), aggregate product-measure
comparison.  It treats all candidates of \(B_k\) at once and does not assert
residual equality, same-index source--target stage equality, paired
outside-\(B_k\) primitive atom equality, or a singleton reduction of \(B_k\).
\end{helper}
\begin{proof}
Apply \(\noderef{SamplingFiniteProductCandidateReplacementStep}\) to the
finite label type \(\alpha\), the parameter \(p=\gamma/\Delta\), the priority
function \(a\mapsto\pi_a\), the one-current-blocker law \(\eta\), and the
independent private-blocker law \(\theta\).  The hypotheses of the cited
finite product theorem are exactly the displayed one-coordinate and
independent-cylinder identities above, together with \(0\le p\le1\) and
\(0\le\pi_a\le1\).  Its conclusion gives
\[
\eta\{\neg(a=1,\ q\in[0,1],\ \pi_b<q\text{ for all }b)\}
\le
\theta\{\exists b:
  \neg(a_b=1,\ q_b\in[0,1],\ \pi_b<q_b)\}.
\]
The two displayed normalization identities identify these two probabilities
with \(\operatorname{ofReal}(P^-_{k,\omega})\) and
\(\operatorname{ofReal}(P^+_{k,\omega})\).  Since
\(P^-_{k,\omega}\) and \(P^+_{k,\omega}\) are nonnegative,
monotonicity and injectivity of \(\operatorname{ofReal}\) on nonnegative
real numbers give \(P^-_{k,\omega}\le P^+_{k,\omega}\).

The remaining displayed mass identities record that these normalized
probabilities are the actual adjacent source and target pieces below the
same fixed noncurrent contexts and paired cells.  They are not used as
independent equality assumptions between residual pieces; they are the
factorization interface which connects the finite product calculation to the
actual replacement step.
\end{proof}
